% !TeX root = ../数学分析教材（伍胜健）笔记.tex
% 本文件由主文档通过 \input 载入；请编译主文档。
\setcounter{chapter}{7} % 保留原书章号。
\chapter{广义积分}
\label{chap:improper-integrals}

% 本章依据“数学分析/textbooks/数学分析（分章原文）/数学分析-8.pdf”正文第 1--26 页整理，不含习题八。
定积分要求积分区间有限且被积函数有界。本章把这两个限制分别放宽：在无穷区间上以截断区间的积分极限定义\emph{无穷积分}（improper integral over an infinite interval），在有限区间上以避开无界点后的积分极限定义\emph{瑕积分}（improper integral with a singularity），二者统称\emph{广义积分}（improper integral）。两类问题都先在正常的闭区间上积分，再研究截断参数的极限；因此，写出广义积分的符号并不表示这个极限已经存在，也不能把定积分的运算法则直接用于尚未证明收敛的表达式。

\section{无穷积分的基本概念与性质}
\label{sec:infinite-integrals-concepts-and-properties}

考虑一个理想化的引力做功问题：设地球质量为 $M$、半径为 $R$，质量固定为 $m$ 的物体沿径向从地表升至高度 $H$，忽略空气阻力、其他天体及地球自转的影响。以 $G$ 表示万有引力常数，则克服引力所做的功为
\[
    W(H)=\int_R^{R+H}\frac{GMm}{r^2}\,\mathrm{d}r
        =GMm\left(\frac1R-\frac1{R+H}\right).
\]
每一个有限的 $H$ 对应普通定积分；物体远离地球时的总功则由 $H\to+\infty$ 的极限给出，等于 $GMm/R$。在同一模型下，把这个极限与初动能 $mv^2/2$ 比较，得到逃逸速度 $v=\sqrt{2GM/R}=\sqrt{2gR}$，其中 $g=GM/R^2$。无穷积分描述的正是这种截断位置不断远移时的累积量。

\begin{definition}[无穷积分及其收敛性]
\label{def:improper-integrals-at-infinity}
设实函数 $f$ 在 $[a,+\infty)$ 上\emph{局部黎曼可积}（locally Riemann integrable），即对每个 $X>a$，$f$ 在 $[a,X]$ 上黎曼可积。若存在有限实数极限
\begin{equation}
    \lim_{X\to+\infty}\int_a^X f(x)\,\mathrm{d}x=I,
    \label{eq:improper-integral-right-limit}
\end{equation}
则称无穷积分 $\int_a^{+\infty}f(x)\,\mathrm{d}x$ \emph{收敛}（convergent），其值定义为 $I$；若无有限极限，则称其\emph{发散}（divergent）。类似地，若 $f$ 在 $(-\infty,b]$ 的每个有限闭子区间上可积，则定义
\[
    \int_{-\infty}^b f(x)\,\mathrm{d}x
    =\lim_{X\to-\infty}\int_X^b f(x)\,\mathrm{d}x,
\]
并以右端有限极限的存在性定义收敛。
\end{definition}

极限等于 $+\infty$ 或 $-\infty$ 仍属于发散；必要时可以用这两个符号说明发散方向，但它们不是收敛积分的值。后文的“局部可积”均指在所讨论区间内每个避开瑕点的有限闭子区间上黎曼可积。

\begin{definition}[整个实轴上的无穷积分]
\label{def:improper-integral-whole-line}
设 $f$ 在实轴的每个有限闭区间上可积。固定 $c\in\mathbb{R}$，只有当 $\int_{-\infty}^c f(x)\,\mathrm{d}x$ 与 $\int_c^{+\infty}f(x)\,\mathrm{d}x$ 分别收敛时，才称 $\int_{-\infty}^{+\infty}f(x)\,\mathrm{d}x$ 收敛，并定义
\begin{equation}
    \int_{-\infty}^{+\infty}f(x)\,\mathrm{d}x
    =\int_{-\infty}^c f(x)\,\mathrm{d}x
     +\int_c^{+\infty}f(x)\,\mathrm{d}x.
    \label{eq:improper-integral-whole-line-split}
\end{equation}
若其中至少一个积分发散，则整个实轴上的积分发散。
\end{definition}

分点 $c$ 的选取不影响敛散性，也不影响收敛时的积分值，因为改变分点只是在两部分之间转移一个有限区间上的定积分。更一般地，对任意 $b>a$，由
\[
    \int_a^X f(x)\,\mathrm{d}x
    =\int_a^b f(x)\,\mathrm{d}x+\int_b^X f(x)\,\mathrm{d}x
    \qquad(X>b)
\]
可知，$\int_a^{+\infty}f$ 与 $\int_b^{+\infty}f$ 同敛散。因此无穷积分的敛散性取决于尾部，而有限初始区间只改变收敛后的数值。若 $f\geq0$，截断积分表示曲线下的面积；若 $f$ 变号，则表示正、负部分面积的代数和，后者可以因抵消而收敛。

\begin{proposition}[收敛积分的线性与原函数计算]
\label{prop:improper-integral-linearity-and-primitives}
若 $\int_a^{+\infty}f$、$\int_a^{+\infty}g$ 均收敛，$\lambda,\mu\in\mathbb{R}$，则
\[
    \int_a^{+\infty}(\lambda f+\mu g)\,\mathrm{d}x
    =\lambda\int_a^{+\infty}f\,\mathrm{d}x
     +\mu\int_a^{+\infty}g\,\mathrm{d}x.
\]
若 $f$ 局部可积且有原函数 $F$，则 $\int_a^{+\infty}f$ 收敛当且仅当 $\lim_{X\to+\infty}F(X)$ 为有限实数，此时
\begin{equation}
    \int_a^{+\infty}f(x)\,\mathrm{d}x
    =\lim_{X\to+\infty}F(X)-F(a).
    \label{eq:improper-newton-leibniz}
\end{equation}
左侧无穷区间上的积分有对应结论；整个实轴上的积分则要求原函数在两端分别有有限极限。
\end{proposition}

\begin{proof}
在每个有限截断区间上，分别使用定积分的线性与牛顿--莱布尼茨公式，再取相应极限即可。线性公式中两项的收敛性保证可以把和的极限拆开；不能把两个发散积分之差当成通常的实数运算。
\end{proof}

\begin{example}
\label{ex:improper-integral-arctangent}
计算 $\int_{-\infty}^{+\infty}\frac{\mathrm{d}x}{1+x^2}$。
\end{example}

\begin{solve}
取分点 $0$，以 $\arctan x$ 为原函数，分别得到
\[
    \int_{-\infty}^0\frac{\mathrm{d}x}{1+x^2}=\frac\pi2,
    \qquad
    \int_0^{+\infty}\frac{\mathrm{d}x}{1+x^2}=\frac\pi2.
\]
两部分均有有限值，故原积分收敛，值为 $\pi$。
\end{solve}

\begin{example}
\label{ex:improper-power-and-log-benchmarks}
讨论幂函数、对数函数构成的基本比较积分。
\end{example}

\begin{solve}
对 $p\in\mathbb{R}$，直接计算截断积分得
\begin{equation}
    \int_1^{+\infty}\frac{\mathrm{d}x}{x^p}
    \begin{cases}
        =\dfrac1{p-1},&p>1,\\
        \text{发散至 }+\infty,&p\leq1.
    \end{cases}
    \label{eq:infinite-power-integral-benchmark}
\end{equation}
这里 $p=1$ 时截断积分为 $\ln X$，$p\neq1$ 时为 $(X^{1-p}-1)/(1-p)$。临界指数 $p=1$ 必须单独检验，不能从 $p\neq1$ 的表达式直接代入。

换元 $t=\ln x$ 后，
\[
    \int_2^{+\infty}\frac{\mathrm{d}x}{x(\ln x)^p}
    =\int_{\ln2}^{+\infty}\frac{\mathrm{d}t}{t^p},
\]
因此仍当且仅当 $p>1$ 时收敛，收敛时的值为 $(\ln2)^{1-p}/(p-1)$。再令 $t=\ln\ln x$，得到
\[
    \int_{e^e}^{+\infty}
        \frac{\mathrm{d}x}{x\ln x(\ln\ln x)^p}
    =\int_1^{+\infty}\frac{\mathrm{d}t}{t^p},
\]
其收敛范围同样为 $p>1$。这些积分说明：幂次达到临界值后，还需由下一层对数因子决定敛散性。
\end{solve}

\begin{proposition}[换元与分部积分]
\label{prop:improper-substitution-and-parts}
设 $f$ 在 $[a,+\infty)$ 上局部可积，$\varphi$ 在 $[\alpha,\beta)$ 上连续可微且严格递增，其中 $\beta$ 可以为 $+\infty$，并满足 $\varphi(\alpha)=a$、$\lim_{t\to\beta-}\varphi(t)=+\infty$。则
\begin{equation}
    \int_a^{+\infty}f(x)\,\mathrm{d}x
    =\int_\alpha^\beta f(\varphi(t))\varphi'(t)\,\mathrm{d}t;
    \label{eq:improper-substitution}
\end{equation}
等式指两边同时收敛或同时发散，收敛时值相等。递减换元应同时调整积分上下限。

若 $u,v\in C^1([a,+\infty))$，且 $L=\lim_{X\to+\infty}u(X)v(X)$ 为有限实数，则 $\int_a^{+\infty}uv'$ 与 $\int_a^{+\infty}u'v$ 同敛散，收敛时
\begin{equation}
    \int_a^{+\infty}u(x)v'(x)\,\mathrm{d}x
    =L-u(a)v(a)-\int_a^{+\infty}u'(x)v(x)\,\mathrm{d}x.
    \label{eq:improper-integration-by-parts}
\end{equation}
这里要求的是乘积 $uv$ 的极限存在，不要求两个因子各有有限极限。
\end{proposition}

\begin{proof}
在 $[\alpha,T]$ 上使用定积分换元公式，右端等于 $\int_a^{\varphi(T)}f$；$T\to\beta-$ 时，$\varphi(T)$ 遍历趋于无穷的截断位置，故两边的极限存在性等价。分部积分则先写出有限区间上的恒等式
\[
    \int_a^X uv'\,\mathrm{d}x
    =u(X)v(X)-u(a)v(a)-\int_a^X u'v\,\mathrm{d}x.
\]
由于边界乘积已有有限极限，右端剩下的积分极限与左端同时存在，从而得到结论。
\end{proof}

\begin{example}
\label{ex:improper-substitution-and-parts-computations}
计算 $\int_1^{+\infty}\frac{\mathrm{d}x}{x\sqrt{1+x^2}}$ 与 $\int_1^{+\infty}\frac{\arctan x}{x^3}\,\mathrm{d}x$。
\end{example}

\begin{solve}
第一个积分令 $t=1/x$，将无穷端点移到 $t=0$，得到
\[
    \int_1^{+\infty}\frac{\mathrm{d}x}{x\sqrt{1+x^2}}
    =\int_0^1\frac{\mathrm{d}t}{\sqrt{1+t^2}}
    =\ln(1+\sqrt2).
\]
换元后的被积函数在 $0$ 处可连续延拓，故右端为普通定积分，这也同时证明原积分收敛。

第二个积分取 $u=\arctan x$、$v=-1/(2x^2)$，则 $uv\to0$，且
\begin{align*}
    \int_1^{+\infty}\frac{\arctan x}{x^3}\,\mathrm{d}x
    &=\frac\pi8+\frac12\int_1^{+\infty}
        \frac{\mathrm{d}x}{x^2(1+x^2)}\\
    &=\frac\pi8+\frac12\int_1^{+\infty}
        \left(\frac1{x^2}-\frac1{1+x^2}\right)\,\mathrm{d}x
      =\frac12.
\end{align*}
右侧各积分均已证明收敛，因此分部积分的使用有依据。
\end{solve}

\begin{example}
\label{ex:exponentially-damped-trigonometric-integrals}
设 $a>0$、$b\in\mathbb{R}$，计算
\[
    I_s=\int_0^{+\infty}e^{-ax}\sin bx\,\mathrm{d}x,
    \qquad
    I_c=\int_0^{+\infty}e^{-ax}\cos bx\,\mathrm{d}x.
\]
\end{example}

\begin{solve}
先验证收敛性：对 $Y>X\geq0$，两个被积函数的尾部积分绝对值均不超过
\[
    \int_X^Y e^{-ax}\,\mathrm{d}x\leq\frac{e^{-aX}}a\longrightarrow0.
\]
因此由函数极限的柯西准则，两个截断积分都有有限极限；下一节的绝对收敛与比较判别法也可直接说明这一点。分别分部积分，边界项因 $e^{-ax}\to0$ 而存在，得
\[
    I_s=\frac ba I_c,\qquad
    I_c=\frac1a-\frac ba I_s.
\]
解这两个线性方程可得
\[
    I_s=\frac{b}{a^2+b^2},\qquad
    I_c=\frac{a}{a^2+b^2}.
    \qedhere
\]
\end{solve}

\begin{definition}[无穷积分的柯西主值]
\label{def:cauchy-principal-value-at-infinity}
设 $f$ 在实轴的每个有限闭区间上可积。若对称截断的有限极限存在，则称
\begin{equation}
    \operatorname{PV}\int_{-\infty}^{+\infty}f(x)\,\mathrm{d}x
    =\lim_{X\to+\infty}\int_{-X}^X f(x)\,\mathrm{d}x
    \label{eq:cauchy-principal-value-infinity}
\end{equation}
为该积分的\emph{柯西主值}（Cauchy principal value）。
\end{definition}

通常收敛要求左右两个极限独立存在，故一定蕴含主值存在且与通常积分相等；主值只考察一种对称截断方式，反向推论不成立。例如 $f(x)=x/(1+x^2)$ 为奇函数，每个对称截断积分都为零，故主值为零，但
\[
    \int_0^X\frac{x}{1+x^2}\,\mathrm{d}x=\frac12\ln(1+X^2)\to+\infty,
\]
所以通常的整个实轴积分发散。\emph[caution]{主值中的对称抵消不能替代两端各自的收敛性检验。}

\section{无穷积分敛散性的判别法}
\label{sec:convergence-tests-for-infinite-integrals}

定义把敛散性归结为 $F(X)=\int_a^X f$ 是否有有限极限，但一般不能求出 $F$ 的初等表达式。判别法的作用是利用被积函数的大小、符号或振荡性质直接控制尾部积分。除特别说明外，本节出现的函数均在 $[a,+\infty)$ 上局部可积；改变有限初始区间不影响结论。

\begin{theorem}[无穷积分的柯西准则]
\label{thm:cauchy-criterion-improper-integrals}
$\int_a^{+\infty}f(x)\,\mathrm{d}x$ 收敛的充分必要条件是：对每个 $\varepsilon>0$，存在 $A>a$，使任意 $v>u>A$ 均满足
\begin{equation}
    \left|\int_u^v f(x)\,\mathrm{d}x\right|<\varepsilon.
    \label{eq:improper-integral-cauchy-tail}
\end{equation}
\end{theorem}

\begin{proof}
令 $F(X)=\int_a^X f(x)\,\mathrm{d}x$，则 $F(v)-F(u)=\int_u^v f$。将函数极限的柯西准则，即定理~\ref{thm:cauchy-criterion-for-function-limits}，用于 $X\to+\infty$，即得结论。
\end{proof}

这里要求\emph{任意两个足够远的截断点之间的积分一致地小}，不是只检验某一列区间或某种配对方式的极限。这个量词要求排除了利用特殊截断掩盖发散的情形。

\begin{definition}[绝对收敛与条件收敛]
\label{def:absolute-and-conditional-improper-integrals}
若 $\int_a^{+\infty}|f(x)|\,\mathrm{d}x$ 收敛，则称 $\int_a^{+\infty}f(x)\,\mathrm{d}x$ \emph{绝对收敛}（absolutely convergent）；若后者收敛而前者发散，则称后者\emph{条件收敛}（conditionally convergent）。
\end{definition}

\begin{theorem}[绝对收敛蕴含收敛]
\label{thm:absolute-improper-integrability-implies-convergence}
若 $\int_a^{+\infty}|f(x)|\,\mathrm{d}x$ 收敛，则 $\int_a^{+\infty}f(x)\,\mathrm{d}x$ 收敛，且
\[
    \left|\int_a^{+\infty}f(x)\,\mathrm{d}x\right|
    \leq\int_a^{+\infty}|f(x)|\,\mathrm{d}x.
\]
\end{theorem}

\begin{proof}
由 $f$ 的局部可积性，$|f|$ 也局部可积。有限区间上的不等式
\[
    \left|\int_u^v f(x)\,\mathrm{d}x\right|
    \leq\int_u^v|f(x)|\,\mathrm{d}x
\]
与 $|f|$ 的柯西准则一起给出 $f$ 的柯西准则，故原积分收敛。再在 $[a,X]$ 上使用同一不等式并令 $X\to+\infty$，得到结论中的估计。
\end{proof}

\begin{theorem}[非负函数积分的有界判据]
\label{thm:nonnegative-improper-integral-boundedness}
若 $f(x)\geq0$，则 $\int_a^{+\infty}f$ 收敛的充分必要条件是其截断积分有共同上界，即存在常数 $C$，使
\[
    \int_a^X f(x)\,\mathrm{d}x\leq C\qquad(X>a).
\]
若无此上界，则截断积分趋于 $+\infty$。
\end{theorem}

\begin{proof}
由定积分对积分区间的可加性，$F(X)=\int_a^X f$ 单调非减。单调函数的极限存在性定理表明，它有有限极限当且仅当有上界；无上界时极限为 $+\infty$。
\end{proof}

\begin{theorem}[比较判别法及其极限形式]
\label{thm:comparison-tests-improper-integrals}
\emph{比较判别法}（comparison test）及其\emph{极限比较形式}（limit comparison test）如下。设 $f,g\geq0$，且存在 $A\geq a$、$C>0$，使 $x\geq A$ 时有 $f(x)\leq Cg(x)$。若 $\int_a^{+\infty}g$ 收敛，则 $\int_a^{+\infty}f$ 收敛；若 $\int_a^{+\infty}f$ 发散，则 $\int_a^{+\infty}g$ 发散。

若进一步有 $g(x)>0$，且 $\lim_{x\to+\infty}f(x)/g(x)=\ell$，则：$0<\ell<+\infty$ 时两积分同敛散；$\ell=0$ 时由 $\int_a^{+\infty}g$ 收敛可推出 $\int_a^{+\infty}f$ 收敛；$\ell=+\infty$ 时由 $\int_a^{+\infty}g$ 发散可推出 $\int_a^{+\infty}f$ 发散。
\end{theorem}

\begin{proof}
若 $\int_a^{+\infty}g$ 收敛，则 $\int_A^Xf\leq C\int_A^Xg$ 有共同上界，由非负函数积分的有界判据得 $\int_A^{+\infty}f$ 收敛，加回有限初始区间即可。发散结论为此命题的逆否命题。若 $0<\ell<+\infty$，充分大时有 $(\ell/2)g\leq f\leq(3\ell/2)g$，故可双向比较；$\ell=0$ 时最终有 $f\leq g$，$\ell=+\infty$ 时最终有 $f\geq g$，分别使用单向比较即可。
\end{proof}

比较的方向可以概括为“大者收敛则小者收敛，小者发散则大者发散”，但前提是被积函数非负。对变号函数，可以先比较绝对值以证明绝对收敛；仅有 $f\sim g$ 而不能控制符号与绝对值时，不能据此照搬上述敛散性结论。极限为零或无穷时也只有定理给出的单向推论。

以例~\ref{ex:improper-power-and-log-benchmarks} 为基准，并利用对数慢于任意正幂增长，可得对 $p,q\in\mathbb{R}$，
\begin{equation}
    \int_e^{+\infty}\frac{\mathrm{d}x}{x^p(\ln x)^q}
    \text{ 收敛}\quad\Longleftrightarrow\quad
    p>1\ \text{或}\ (p=1\text{ 且 }q>1).
    \label{eq:power-log-infinite-integral-test}
\end{equation}
具体地，$p>1$ 时选 $1<r<p$，有 $x^{-p}(\ln x)^{-q}=o(x^{-r})$；$p<1$ 时选 $p<r<1$，则该被积函数与 $x^{-r}$ 的比值趋于 $+\infty$；$p=1$ 时令 $t=\ln x$。同理，对任意多项式 $P$ 和 $a>0$，$|P(x)|e^{-ax}\leq Ce^{-ax/2}$ 在充分大时成立，所以 $\int_0^{+\infty}P(x)e^{-ax}\,\mathrm{d}x$ 绝对收敛；这也给出例~\ref{ex:exponentially-damped-trigonometric-integrals} 中两个积分的另一种收敛性证明。

\begin{example}
\label{ex:rational-infinite-integral-degree-test}
设 $n$ 为正整数，讨论 $\int_0^{+\infty}\frac{x^2}{x^n+x^2+1}\,\mathrm{d}x$ 的敛散性，并说明有理函数的一般判据。
\end{example}

\begin{solve}
被积函数在每个有限闭区间上连续。$n>2$ 时，它在无穷远处等价于 $x^{2-n}$，故当且仅当 $n-2>1$，即 $n>3$ 时收敛。$n=1$、$n=2$ 时，被积函数分别趋于 $1$、$1/2$，故均发散；这两种情形的分母最高次项不同，不能直接套用 $n>2$ 时的等价式。

一般地，设非零实多项式 $P,Q$ 的次数分别为 $m,n$，首项系数分别为 $p_m,q_n$，且 $Q$ 在 $[a,+\infty)$ 上无零点，则
\[
    \frac{P(x)}{Q(x)}\sim\frac{p_m}{q_n}x^{m-n}
    \qquad(x\to+\infty).
\]
有理函数在充分远处符号固定，故其积分收敛等价于绝对收敛，再由幂积分判据可知，这当且仅当 $n-m\geq2$。若有限区间中还有未消去的极点，则必须另行检验，次数判据只处理无穷端。
\end{solve}

\begin{example}
\label{ex:absolute-comparison-and-log-cancellation}
证明 $\int_1^{+\infty}\frac{\cos x\sin(1/x)}x\,\mathrm{d}x$ 绝对收敛，并讨论 $p>0$ 时 $\int_1^{+\infty}[x^{-p}-\ln(1+x^{-p})]\,\mathrm{d}x$ 的敛散性。
\end{example}

\begin{solve}
由 $|\sin t|\leq|t|$，第一问的被积函数绝对值不超过 $x^{-2}$，故绝对收敛。第二问利用 $t-\ln(1+t)=t^2/2+o(t^2)$，令 $t=x^{-p}$，得到
\[
    x^{-p}-\ln(1+x^{-p})\sim\frac1{2x^{2p}}.
\]
被积函数非负，故极限比较合法，原积分当且仅当 $2p>1$，即 $p>1/2$ 时收敛。这里首阶项相消，必须展开到二阶；若只把 $\ln(1+x^{-p})$ 换成 $x^{-p}$，便会丢失决定敛散性的主项。
\end{solve}

\begin{example}
\label{ex:continuous-integrable-spikes}
构造连续正函数 $f$，使 $\int_1^{+\infty}f(x)\,\mathrm{d}x$ 收敛，但 $f(x)$ 在无穷远处不趋于零，且无界。
\end{example}

\begin{solve}
用越来越窄的三角形尖峰控制总面积。记 $r_+=\max\{r,0\}$，在每个整数 $n\geq2$ 附近放置高为 $2^n$、半底宽为 $3^{-n}$ 的尖峰，定义
\[
    f(x)=\frac1{x^2}
       +\sum_{n=2}^{\infty}2^n\bigl(1-3^n|x-n|\bigr)_+,
    \qquad x\geq1.
\]
各尖峰支撑区间 $[n-3^{-n},n+3^{-n}]$ 互不相交，且任意有限闭区间只遇到有限个尖峰，所以这里的和在局部实际为有限和，$f$ 连续且严格为正。第 $n$ 个三角形的面积为 $(2/3)^n$，故对任意 $X>1$，
\[
    \int_1^X f(x)\,\mathrm{d}x
    \leq1+\sum_{n=2}^{\infty}\left(\frac23\right)^n
    =\frac73.
\]
由非负函数积分的有界判据，广义积分收敛；但 $f(n)=n^{-2}+2^n\to+\infty$，所以函数值并不趋于零。这一论证只用了有限区间上有限多个尖峰的面积估计，没有交换无穷求和与积分。
\end{solve}

\emph[caution]{积分收敛控制的是尾部累积量，一般不强迫被积函数逐点趋于零，即使函数连续且非负也不例外。}反过来，$f(x)\to0$ 也不足以保证收敛，$f(x)=1/x$ 即为反例。若另外知道 $f$ 在尾部单调，则积分收敛必推出 $f(x)\to0$：单调性保证扩充实数极限存在，若极限非零，则尾部最终符号固定且绝对值有正下界，其积分不可能有有限极限。

\begin{theorem}[狄利克雷判别法]
\label{thm:dirichlet-test-infinite-integrals}
设 $f$ 局部可积，其累积积分一致有界，即存在 $M>0$，使
\[
    \left|\int_a^X f(x)\,\mathrm{d}x\right|\leq M
    \qquad(X\geq a).
\]
若 $g$ 在 $[a,+\infty)$ 上单调且 $\lim_{x\to+\infty}g(x)=0$，则 $\int_a^{+\infty}f(x)g(x)\,\mathrm{d}x$ 收敛。这称为\emph{狄利克雷判别法}（Dirichlet's test）。
\end{theorem}

\begin{proof}
记 $F(X)=\int_a^X f$，则任意 $v>u\geq a$ 有 $|\int_u^v f|=|F(v)-F(u)|\leq2M$。在有限区间 $[u,v]$ 上，$g$ 单调，故第二积分中值定理给出某个 $\xi\in[u,v]$，使
\[
    \int_u^v fg\,\mathrm{d}x
    =g(u)\int_u^\xi f\,\mathrm{d}x
     +g(v)\int_\xi^v f\,\mathrm{d}x.
\]
因此
\[
    \left|\int_u^v fg\,\mathrm{d}x\right|
    \leq2M\bigl(|g(u)|+|g(v)|\bigr).
\]
当 $u,v$ 都足够大时，右端可任意小；由柯西准则，乘积积分收敛。单调函数在有限区间上可积，故 $fg$ 的局部可积性也有保证。
\end{proof}

\begin{theorem}[阿贝尔判别法]
\label{thm:abel-test-infinite-integrals}
若 $\int_a^{+\infty}f(x)\,\mathrm{d}x$ 收敛，而 $g$ 在 $[a,+\infty)$ 上单调有界，则 $\int_a^{+\infty}f(x)g(x)\,\mathrm{d}x$ 收敛。这称为\emph{阿贝尔判别法}（Abel's test）。
\end{theorem}

\begin{proof}
由 $g$ 单调有界，存在有限极限 $\ell=\lim_{x\to+\infty}g(x)$。又因 $\int_a^Xf$ 收敛于有限值，其在尾部有界，在每个有限初始区间上也有界，故可将狄利克雷判别法用于 $f$ 与 $g-\ell$，得到 $\int_a^{+\infty}f(g-\ell)$ 收敛。再加上已收敛的 $\ell\int_a^{+\infty}f$ 即得结论。
\end{proof}

两种判别法的分工是：狄利克雷判别法只要求一个因子的累积积分有界，但要求另一个因子单调趋于零；阿贝尔判别法要求第一个因子的积分已经收敛，另一个因子只需单调有界。例如，在下例证明 $\int_1^{+\infty}\sin x/x\,\mathrm{d}x$ 收敛后，再乘以单调有界的 $\arctan x$，即可用阿贝尔判别法判断所得积分收敛。两种判别法通常给出收敛性，绝对收敛仍需另行检验。

\begin{example}
\label{ex:sinc-integral-conditional-convergence}
证明 $\int_{-\infty}^{+\infty}\frac{\sin x}{x}\,\mathrm{d}x$ 条件收敛，其中在 $x=0$ 处把被积函数连续延拓为 $1$。
\end{example}

\begin{solve}
函数在有限区间上连续，只需讨论无穷尾部。在 $[1,+\infty)$ 上，$\int_1^X\sin x\,\mathrm{d}x=\cos1-\cos X$ 一致有界，而 $1/x$ 单调趋于零，故由狄利克雷判别法，$\int_1^{+\infty}\sin x/x\,\mathrm{d}x$ 收敛。被积函数为偶函数，负半轴上的积分也收敛，因而整个实轴上的通常积分收敛。

另一方面，对 $X>1$，利用 $|\sin x|\geq\sin^2x$，有
\[
    \int_1^X\frac{|\sin x|}{x}\,\mathrm{d}x
    \geq\int_1^X\frac{\sin^2x}{x}\,\mathrm{d}x
    =\frac12\ln X-\frac12\int_1^X\frac{\cos2x}{x}\,\mathrm{d}x.
\]
最后一个振荡积分由狄利克雷判别法收敛，右端因 $\ln X\to+\infty$ 而趋于 $+\infty$。所以绝对值积分发散，原积分条件收敛。这里先判定收敛而未求其值，说明判别法可以处理没有初等原函数的积分。
\end{solve}

\begin{example}
\label{ex:oscillatory-log-improper-integral}
固定 $A>1$，讨论参数 $p>0$ 时
\[
    I_p=\int_A^{+\infty}\ln\left(1+\frac{\sin x}{x^p}\right)\,\mathrm{d}x
\]
的收敛性与绝对收敛性。
\end{example}

\begin{solve}
由于 $|\sin x|/x^p<1$，对数在整个积分区间上有定义。令 $u(x)=\sin x/x^p$，并设
\[
    d(x)=u(x)-\ln(1+u(x)).
\]
由 $\ln(1+u)\leq u$ 以及泰勒展开，在充分大的 $x$ 上有
\[
    d(x)=\frac{u(x)^2}{2}\bigl(1+O(u(x))\bigr),\qquad
    \frac14\frac{\sin^2x}{x^{2p}}
    \leq d(x)\leq\frac34\frac{\sin^2x}{x^{2p}}.
\]
在 $u=0$ 的点两侧均为零，上述不等式仍成立，因而不必在那里作除法。对任意 $p>0$，$\int_A^{+\infty}u(x)\,\mathrm{d}x$ 均由狄利克雷判别法收敛，而
\[
    \int_A^X\frac{\sin^2x}{x^{2p}}\,\mathrm{d}x
    =\frac12\int_A^X\frac{\mathrm{d}x}{x^{2p}}
     -\frac12\int_A^X\frac{\cos2x}{x^{2p}}\,\mathrm{d}x.
\]
振荡项对全部 $p>0$ 收敛，故非负积分 $\int_A^{+\infty}d$ 当且仅当 $2p>1$ 时收敛。于是 $0<p\leq1/2$ 时，$I_p$ 的截断积分趋于 $-\infty$；$p>1/2$ 时，$I_p$ 收敛。

为区分绝对收敛，注意 $|\ln(1+u)|$ 在 $u\to0$ 时与 $|u|$ 可双向作常数倍比较。$p>1$ 时，$|u(x)|\leq x^{-p}$，故绝对收敛。$0<p\leq1$ 时，在每个区间 $[n\pi+\pi/6,n\pi+5\pi/6]$ 上均有 $|\sin x|\geq1/2$，所以
\[
    \int_{n\pi}^{(n+1)\pi}\frac{|\sin x|}{x^p}\,\mathrm{d}x
    \geq\frac{\pi}{3((n+1)\pi)^p}.
\]
对充分大的连续整数 $n$ 求和，并与 $\int t^{-p}\,\mathrm{d}t$ 比较，得到绝对值积分发散。因此最终分类为：$p>1$ 时绝对收敛，$1/2<p\leq1$ 时条件收敛，$0<p\leq1/2$ 时发散。

尽管在 $\sin x\neq0$ 且 $x\to+\infty$ 时有 $\ln(1+u(x))/u(x)\to1$，当 $p\leq1/2$ 时它的积分与 $u$ 的积分却不同敛散。首阶振荡项的抵消使二阶的非负累积量成为关键，这正是变号情形不能直接套用等价比较的原因。
\end{solve}

\section{瑕积分}
\label{sec:singular-integrals}

\subsection{瑕积分的概念}
\label{subsec:concept-of-singular-integrals}

有些函数只在有限区间的端点或少数内点附近无界，在避开这些点的闭区间上仍可积。无界性使它不满足普通黎曼积分的必要条件，却不排除截断积分有有限极限。若函数在某点某一侧的每个充分小去心邻域内均无界，就称该点为该侧的\emph{瑕点}（singular point）；瑕积分考察的正是逼近瑕点时能否得到确定的有限累积量。

\begin{definition}[端点处的瑕积分]
\label{def:singular-integrals-at-endpoints}
设 $f$ 在 $(a,b]$ 上有定义，在每个 $[a+\varepsilon,b]$ 上可积，其中 $0<\varepsilon<b-a$。若有限极限存在，则定义
\begin{equation}
    \int_a^b f(x)\,\mathrm{d}x
    =\lim_{\varepsilon\downarrow0}
        \int_{a+\varepsilon}^b f(x)\,\mathrm{d}x,
    \label{eq:singular-integral-left-endpoint}
\end{equation}
并称相应瑕积分收敛；否则称其发散。若无界端点为 $b$，则相应定义为
\begin{equation}
    \int_a^b f(x)\,\mathrm{d}x
    =\lim_{\varepsilon\downarrow0}
        \int_a^{b-\varepsilon} f(x)\,\mathrm{d}x.
    \label{eq:singular-integral-right-endpoint}
\end{equation}
\end{definition}

瑕点处是否定义函数值及如何赋值不影响上述极限。另一方面，单纯在端点缺少定义却可有界延拓的函数，可能直接补成普通可积函数，并不一定构成真正的无界瑕点。

若两端都是瑕点，应在 $c\in(a,b)$ 处分成 $\int_a^c f+\int_c^b f$，要求两部分分别收敛。若内部点 $c$ 为瑕点，则同样定义
\[
    \int_a^b f(x)\,\mathrm{d}x
    =\lim_{\varepsilon\downarrow0}\int_a^{c-\varepsilon}f(x)\,\mathrm{d}x
     +\lim_{\eta\downarrow0}\int_{c+\eta}^b f(x)\,\mathrm{d}x,
\]
两个极限必须独立存在；有限多个瑕点时依此逐段处理。若另有无穷端点，也须分开检验每个异常端点，不能让一个瑕点附近的发散与另一个端点的发散互相抵消。

对内部瑕点，也可另外定义\emph{柯西主值}
\[
    \operatorname{PV}\int_a^b f(x)\,\mathrm{d}x
    =\lim_{\varepsilon\downarrow0}
       \left(\int_a^{c-\varepsilon}f(x)\,\mathrm{d}x
             +\int_{c+\varepsilon}^b f(x)\,\mathrm{d}x\right),
\]
这里假定其余端点正常。它与定义~\ref{def:cauchy-principal-value-at-infinity} 一样，只考察特定的同步截断。例如 $\operatorname{PV}\int_{-1}^1\mathrm{d}x/x=0$，但两侧积分分别发散，所以通常的瑕积分不存在。

若 $f$ 在 $(a,b]$ 上局部可积且有原函数 $F$，则逐个截断后使用牛顿--莱布尼茨公式得到
\begin{equation}
    \int_a^b f(x)\,\mathrm{d}x
    =F(b)-\lim_{x\to a+}F(x),
    \label{eq:singular-newton-leibniz}
\end{equation}
其中右端要求原函数的单侧极限为有限实数；右端点为瑕点时作相应修改。原函数有界本身并不足以保证该极限存在，关键仍是极限而非有界性。

\begin{example}
\label{ex:unbounded-derivative-singular-integral}
设 $F(0)=0$，$F(x)=x^2\sin(1/x^2)$（$0<|x|\leq1$），讨论其导函数的积分。
\end{example}

\begin{solve}
由导数定义有 $F'(0)=\lim_{h\to0}h\sin(1/h^2)=0$，而对 $x>0$，
\[
    f(x)=F'(x)=2x\sin\frac1{x^2}-\frac2x\cos\frac1{x^2}.
\]
取 $x_n=(2\pi n)^{-1/2}$，便有 $f(x_n)=-2/x_n\to-\infty$，故 $f$ 在零点附近无界，不能在 $[0,1]$ 上作普通黎曼积分。但它在每个 $[\varepsilon,1]$ 上连续，且 $F(\varepsilon)\to0$，所以
\[
    \int_0^1 f(x)\,\mathrm{d}x
    =\lim_{\varepsilon\downarrow0}
        \bigl(F(1)-F(\varepsilon)\bigr)=\sin1.
\]
这说明有原函数并不保证普通黎曼可积，却可以借原函数的端点极限计算瑕积分。
\end{solve}

\begin{example}
\label{ex:singular-arccosine-and-logarithm}
计算 $\int_{-1}^1\frac{\arccos x}{\sqrt{1-x^2}}\,\mathrm{d}x$ 与 $\int_0^1\ln x\,\mathrm{d}x$。
\end{example}

\begin{solve}
第一问的被积函数在 $x\to1-$ 时极限为 $1$，可作连续延拓；在 $x=-1$ 附近才有无界瑕点。以 $-\tfrac12(\arccos x)^2$ 为原函数，逐端取极限可得
\[
    \int_{-1}^1\frac{\arccos x}{\sqrt{1-x^2}}\,\mathrm{d}x
    =0-\left(-\frac{\pi^2}{2}\right)=\frac{\pi^2}{2}.
\]
第二问在 $x=0$ 附近无界。先在 $[\varepsilon,1]$ 上分部积分，再利用 $\varepsilon\ln\varepsilon\to0$，得到
\[
    \int_0^1\ln x\,\mathrm{d}x
    =\lim_{\varepsilon\downarrow0}
        [x\ln x-x]_\varepsilon^1=-1.
    \qedhere
\]
\end{solve}

\begin{example}
\label{ex:singular-power-integral-benchmark}
设 $a<b$、$p\in\mathbb{R}$，讨论 $\int_a^b(b-x)^{-p}\,\mathrm{d}x$ 的敛散性。
\end{example}

\begin{solve}
若 $p=1$，截断积分为 $\ln(b-a)-\ln\varepsilon\to+\infty$；若 $p\neq1$，则
\[
    \int_a^{b-\varepsilon}\frac{\mathrm{d}x}{(b-x)^p}
    =\frac{(b-a)^{1-p}-\varepsilon^{1-p}}{1-p}.
\]
因此
\begin{equation}
    \int_a^b\frac{\mathrm{d}x}{(b-x)^p}
    \begin{cases}
        =\dfrac{(b-a)^{1-p}}{1-p},&p<1,\\
        \text{发散至 }+\infty,&p\geq1.
    \end{cases}
    \label{eq:singular-power-integral-benchmark}
\end{equation}
$p\leq0$ 时被积函数在端点可连续延拓，实际为普通定积分；$0<p<1$ 时虽无界却可积。左端点的 $\int_a^b(x-a)^{-p}\,\mathrm{d}x$ 有相同结论。与式~\eqref{eq:infinite-power-integral-benchmark} 对照，无穷远处要求 $p>1$，有限瑕点附近要求 $p<1$，不能混用。
\end{solve}

两类广义积分可以通过换元相互转化。对右端点 $b$，令 $t=1/(b-x)$，先在截断区间上换元再取极限，有
\begin{equation}
    \int_a^b f(x)\,\mathrm{d}x
    =\int_{1/(b-a)}^{+\infty}
        f\left(b-\frac1t\right)\frac{\mathrm{d}t}{t^2}.
    \label{eq:singular-to-infinite-integral}
\end{equation}
相应地，左端点的瑕积分可写成 $\int_{1/(b-a)}^{+\infty}f(a+1/t)t^{-2}\,\mathrm{d}t$。反过来，对 $a>0$，令 $t=1/x$ 可得
\[
    \int_a^{+\infty}f(x)\,\mathrm{d}x
    =\int_0^{1/a}f(1/t)\frac{\mathrm{d}t}{t^2}.
\]
这些等式均表示两端极限的等价性；右端可能是瑕积分，也可能在补定义后成为普通定积分。尤其对幂函数，式~\eqref{eq:singular-to-infinite-integral} 把 $(b-x)^{-p}$ 化为 $t^{p-2}$，无穷积分条件 $2-p>1$ 正好等价于瑕积分条件 $p<1$。瑕积分的换元和分部积分也应遵循同一原则：先在避开瑕点的闭区间上作普通运算，再核对换元后的端点极限；分部积分还需检验边界乘积的有限极限，不能直接在瑕点处代入无定义的量。

\subsection{瑕积分敛散性的判别法}
\label{subsec:convergence-tests-for-singular-integrals}

以下以 $b$ 为唯一可能的无界端点，设所有函数均在 $[a,b)$ 上局部可积。左端点的结论可由反射换元 $x=a+b-t$ 得到。若不便直接求原函数或换元后的表达式更复杂，可以在原区间上使用与无穷积分平行的尾部判据。

\begin{theorem}[瑕积分的柯西准则与绝对收敛]
\label{thm:cauchy-and-absolute-singular-integrals}
$\int_a^b f(x)\,\mathrm{d}x$ 收敛的充分必要条件是：对任意 $\varepsilon>0$，存在 $\delta\in(0,b-a)$，使任意 $b-\delta<u<v<b$ 均有
\begin{equation}
    \left|\int_u^v f(x)\,\mathrm{d}x\right|<\varepsilon.
    \label{eq:singular-integral-cauchy-tail}
\end{equation}
若 $\int_a^b|f(x)|\,\mathrm{d}x$ 收敛，称原瑕积分绝对收敛，此时原积分必收敛；若原积分收敛而绝对值积分发散，则称其条件收敛。
\end{theorem}

\begin{proof}
令 $F(t)=\int_a^t f$，其在 $t\to b-$ 时的函数极限柯西准则恰为式~\eqref{eq:singular-integral-cauchy-tail}。绝对收敛蕴含收敛仍由 $|\int_u^v f|\leq\int_u^v|f|$ 与柯西准则得到。
\end{proof}

\begin{theorem}[瑕积分的比较判别法]
\label{thm:comparison-tests-singular-integrals}
设 $f,g\geq0$。若存在 $\delta_0\in(0,b-a)$、$C>0$，使 $b-\delta_0<x<b$ 时有 $f(x)\leq Cg(x)$，则 $\int_a^b g$ 收敛蕴含 $\int_a^b f$ 收敛，$\int_a^b f$ 发散蕴含 $\int_a^b g$ 发散。

若 $g>0$ 且 $\lim_{x\to b-}f(x)/g(x)=\ell$，则 $0<\ell<+\infty$ 时两积分同敛散；$\ell=0$ 时由 $\int_a^b g$ 收敛可推得 $\int_a^b f$ 收敛；$\ell=+\infty$ 时由 $\int_a^b g$ 发散可推得 $\int_a^b f$ 发散。
\end{theorem}

\begin{proof}
对非负函数，$\int_a^t f$ 随 $t\uparrow b$ 单调非减，有限极限的存在仍等价于这些截断积分有共同上界。因此邻近端点的常数倍比较可转化为截断积分的比较，而远离端点的部分只是有限定积分。极限形式通过在端点附近建立常数倍不等式推出，与定理~\ref{thm:comparison-tests-improper-integrals} 的证明相同。
\end{proof}

最基本的比较对象是 $(b-x)^{-p}$，其收敛阈值见式~\eqref{eq:singular-power-integral-benchmark}。例如 $f\geq0$ 且 $f(x)\sim C(b-x)^{-p}$、$C>0$ 时，积分当且仅当 $p<1$ 时收敛。比较前应先确定全部瑕点，并在每个瑕点附近分别找主项；函数在某一端的等价式不能说明另一端的行为。

\begin{example}
\label{ex:two-endpoint-singular-power-integral}
讨论 $\int_{-1}^1(1-x^2)^{-p}\,\mathrm{d}x$ 的敛散性，其中 $p\in\mathbb{R}$。
\end{example}

\begin{solve}
在 $x\to1-$ 与 $x\to-1+$ 时分别有
\[
    (1-x^2)^{-p}\sim2^{-p}(1-x)^{-p},\qquad
    (1-x^2)^{-p}\sim2^{-p}(1+x)^{-p}.
\]
在 $0$ 处分开后，两部分由同一个幂函数判据控制，故原积分当且仅当 $p<1$ 时收敛；$p\geq1$ 时两端均发散。$p\leq0$ 时端点可作有界连续延拓，属于普通可积情形，不能不加区分地说所有参数下都有两个无界瑕点。
\end{solve}

\begin{theorem}[瑕积分的狄利克雷与阿贝尔判别法]
\label{thm:dirichlet-abel-singular-integrals}
若存在 $M>0$，使任意 $a\leq t<b$ 均有 $|\int_a^t f(x)\,\mathrm{d}x|\leq M$，而 $g$ 在 $[a,b)$ 上单调且 $\lim_{x\to b-}g(x)=0$，则 $\int_a^b f(x)g(x)\,\mathrm{d}x$ 收敛；这是瑕积分的狄利克雷判别法。若 $\int_a^b f(x)\,\mathrm{d}x$ 已收敛，而 $g$ 在 $[a,b)$ 上单调有界，则乘积积分也收敛；这是瑕积分的阿贝尔判别法。
\end{theorem}

\begin{proof}
第一种情形下，对 $a\leq u<v<b$，有限区间上的第二积分中值定理与累积积分有界性给出
\[
    \left|\int_u^v fg\,\mathrm{d}x\right|
    \leq2M\bigl(|g(u)|+|g(v)|\bigr).
\]
当 $u,v\to b-$ 时右端一致趋于零，故由瑕积分的柯西准则收敛。第二种情形下，$g$ 有有限端点极限 $\ell$，而 $\int_a^t f$ 有界；先把第一种情形用于 $f(g-\ell)$，再加回已收敛的 $\ell\int_a^b f$ 即可。
\end{proof}

\begin{example}
\label{ex:mixed-logarithmic-improper-integral}
讨论含无穷端点与有限异常端点的积分
\[
    \int_0^{+\infty}\frac{\ln(1+x)}{x^\alpha}\,\mathrm{d}x,
    \qquad\alpha\in\mathbb{R}.
\]
\end{example}

\begin{solve}
在 $1$ 处分成 $\int_0^1$ 与 $\int_1^{+\infty}$。当 $x\to0+$ 时，$\ln(1+x)\sim x$，故
\[
    \frac{\ln(1+x)}{x^\alpha}\sim x^{1-\alpha}.
\]
由左端点的幂函数判据，零点附近的积分当且仅当 $\alpha<2$ 时收敛；其中 $\alpha\leq1$ 时可连续延拓，$1<\alpha<2$ 时虽无界却可积。当 $x\to+\infty$ 时，$\ln(1+x)\sim\ln x$，由式~\eqref{eq:power-log-infinite-integral-test}，无穷尾部当且仅当 $\alpha>1$ 时收敛。两项必须同时满足，故原积分收敛的充分必要条件为
\[
    1<\alpha<2.
\]
临界值 $\alpha=1$ 在无穷端发散，$\alpha=2$ 在零点端发散；将两类端点混在同一个截断过程中计算，可能掩盖其中一个条件。
\end{solve}

\section*{习题八}
\addcontentsline{toc}{section}{习题八}
